\section{Saturated graphs and arithmetic}\label{un:arithmetic}

We now use bounded numerical comprehension to recover the graph
constructions needed for arithmetic. The all-representatives coding and
the canonical history construction are inherited from
draft~\cite{draftC2}; their existence premises now follow from the unary
theory. No graph object is used as a free parameter of a filter matrix.

\subsection{Singletons, pairs, and ordered-pair codes}

Define the following fixed formulas:
\begin{align*}
 \operatorname{SSing}(h,x)&:\Longleftrightarrow
             \forall z\,(z\in h\leftrightarrow z\Eq x),\\
 \operatorname{SPair}(d,x,y)&:\Longleftrightarrow
             \forall z\,(z\in d\leftrightarrow(z\Eq x\lor z\Eq y)),\\
 \operatorname{OP}(p,x,y)&:\Longleftrightarrow
       \N(x)\land\N(y)\land
       \forall u\,[u\in p\leftrightarrow
          (\operatorname{SSing}(u,x)\lor\operatorname{SPair}(u,x,y))].
\end{align*}
The last construction collects every witness for the two Kuratowski
component profiles, not an arbitrarily chosen representative of each.

\begin{lemma}[Saturated numerical codes]\label{un:codes}
Every number has an $\operatorname{SSing}$ witness, and every pair of
numbers has $\operatorname{SPair}$ and $\operatorname{OP}$ witnesses.
These formulas are invariant under $\Eq$ in all their arguments, and
their outputs are saturated and unique modulo $\Eq$. An
$\operatorname{OP}$ code determines both numerical coordinates uniquely
modulo $\Eq$. If $x,y\in n$, $\N(n)$, and $\Pow(n,W)$, then every
$\operatorname{OP}(p,x,y)$ satisfies $p\subseteq W$.
\end{lemma}
\begin{proof}
For a numerical $x$, use Theorem~\ref{un:bnc} on the target
$z\Eq x$, supported at $x$. For numerical $x,y$, choose a common
number $n$ containing both, using Corollary~\ref{un:bootstrap}.
The target $z\Eq x\lor z\Eq y$ is invariant in its numerical
parameters and is supported at $n$: its selected objects are
coextensional with members of the transitive object $n$.
This gives $\operatorname{SPair}$.

If $x,y\in n$, saturation of $n$ shows that every singleton or pair
witness for those coordinates is a subset of $n$. Therefore the target
\[
             \operatorname{SSing}(u,x)\lor\operatorname{SPair}(u,x,y)
\]
is supported at $n$. It is invariant in the numerical parameters $x,y$.
Another application of Theorem~\ref{un:bnc} gives the
$\operatorname{OP}$ witness.

Each output profile is unchanged when its inputs are changed within
their coextensionality classes. Its membership test is invariant in the
tested object as well, which proves saturation. The displayed exact
profiles give output uniqueness modulo $\Eq$.

For decoding, an $\operatorname{OP}$ witness satisfies
\begin{align}
 \forall u\,(u\in p\to z\in u)
       &\quad\longleftrightarrow\quad z\Eq x,\label{un:decode-first}\\
 \exists u\,(u\in p\land z\in u)
       &\quad\longleftrightarrow\quad(z\Eq x\lor z\Eq y).
                                                        \label{un:decode-both}
\end{align}
A singleton witness exists and belongs to $p$, so it forces the
right-hand condition of~\eqref{un:decode-first}; every possible
component contains that entire class, proving the converse.
A pair witness exists and belongs to $p$, proving both directions
of~\eqref{un:decode-both}. The first formula determines $[x]$ and
the second determines $\{[x],[y]\}$; together they determine $[y]$,
including the case $[x]=[y]$.

Finally, every $u\in p$ is a singleton or pair witness and hence
satisfies $u\subseteq n$. The power identity puts $u$ in $W$.
This is precisely $p\subseteq W$.
\end{proof}

\subsection{Comprehension for numerically parameterized graphs}

For any object $G$ put
\[
       V_G(x,y):=\exists p\,(p\in G\land\operatorname{OP}(p,x,y)).
\]
The numerical coordinates of this relation are invariant under $\Eq$,
even if $G$ is an arbitrary ambient object. Replacing $G$ by a
coextensional object preserves the relation as well.

\begin{theorem}[Numerical graph comprehension]\label{un:graphs}
Let $R(x,y,n,\bar r)$ be any fixed formula whose only free variables
are the displayed numerical ones, and suppose its invariance under
$\Eq$ in all those arguments has been proved. Then for every numerical
$n,\bar r$ there exists a saturated $G$ such that
\begin{equation}\label{un:graph-profile}
 p\in G\quad\longleftrightarrow\quad
 \exists x\exists y\,
 [x\in n\land y\in n\land\operatorname{OP}(p,x,y)
                                  \land R(x,y,n,\bar r)].
\end{equation}
Its decoded relation satisfies
\begin{equation}\label{un:graph-relation}
 V_G(x,y)\quad\longleftrightarrow\quad
              x\in n\land y\in n\land R(x,y,n,\bar r).
\end{equation}
The parameters $\bar r$ need not themselves belong to $n$.
\end{theorem}
\begin{proof}
The right side of~\eqref{un:graph-profile}, viewed as a predicate of
$p$ with numerical parameter tuple $(n,\bar r)$, is invariant in those
parameters. Coextensionality of $n$ preserves its member profile, and
the other assertions are invariant by hypothesis and
Lemma~\ref{un:codes}. Choose an exact power $W$ of the number $n$,
using Corollary~\ref{un:bootstrap}. Every selected code is a subset of
$W$, by Lemma~\ref{un:codes}. Thus Theorem~\ref{un:bnc} supplies
the exact output $G$. The membership test depends only on the profile
of $p$, so $G$ is saturated.

For the reverse implication in~\eqref{un:graph-relation}, coordinates
belonging to $n$ are numerical by heredity, and an ordered-pair code
exists;~\eqref{un:graph-profile} puts it into $G$. Conversely, a code
witnessing $V_G(x,y)$ is selected by some coordinates $x',y'\in n$
satisfying $R$. Unique decoding gives $x\Eq x'$ and $y\Eq y'$.
Saturation puts $x,y$ in $n$, and invariance of $R$ proves the desired
right side.
\end{proof}

Only a power of a numerical source is needed here. The source-free
proof uses neither iterated powers nor a formula recovering a source
from its power. The bound and every parameter are bound inside the
least-failure matrix of Theorem~\ref{un:bnc}.

\subsection{Deterministic histories with explicit value bounds}

Fix a formula $T(u,v;\bar t)$ with only the displayed free variables.
Its parameters, input, and output are numerical. Assume it is invariant
under $\Eq$ and defines a total function on numerical classes:
\[
\begin{split}
 &\forall u\in\N\;\exists v\in\N\;T(u,v;\bar t),\\
 &u\Eq u'\land T(u,v;\bar t)\land T(u',v';\bar t)
                                      \quad\to\quad v\Eq v',
\end{split}
\]
where the numerical guards in the second line are understood.
These assumptions will be proved for the particular transitions used
to define addition and multiplication.

For numerical $a,b,c,d,\bar t$, define
$H_T(f;a,b,c,d;\bar t)$ by the following finite first-order conditions:
\begin{enumerate}
\item Every member of $f$ is an $\operatorname{OP}$ code with numerical
      coordinates $i\preccurlyeq b$, $y\preccurlyeq d$.
\item Each numerical $i\preccurlyeq b$ has a decoded value
      $y\preccurlyeq d$, and any two values at the same index are
      coextensional.
\item For every empty $e$, $V_f(e,a)$ holds, and $V_f(b,c)$ holds.
\item If $\Succ(i,j)$, $j\preccurlyeq b$, $V_f(i,u)$ and $V_f(j,v)$,
      with numerical arguments, then $T(u,v;\bar t)$.
\end{enumerate}
The literal formula also includes $\N$-guards on every free numerical
argument. The bound $d$ is separate from the terminal value $c$; no
monotonicity of $T$ is assumed. The definition is invariant in every
numerical argument, directly from the invariance of comparison,
$\operatorname{OP}$, and $T$. This fact is established before the
existence or uniqueness of histories.

\begin{lemma}[Pointwise comparison before history existence]
\label{un:history-unique}
Two $T$-histories with coextensional initial values have coextensional
values at every index common to their domains. Their lengths and
numerical value bounds may differ.
\end{lemma}
\begin{proof}
Keep the two histories and all their parameters fixed. Apply full
representative induction, Theorem~\ref{un:induction}, to the assertion
that their values agree at $i$ whenever $i$ belongs to both domains.
At an empty index, their initial clauses and the functionality of each
graph identify its values with its initial class.

For an existing successor $\Succ(i,j)$, if $j$ belongs to both domains,
so does $i$, because $i\in j$ and those numerical segments are initial.
Take the two values at $i$ from graph totality. They are coextensional
by the induction hypothesis. Each value at $j$ is a $T$-output of the
corresponding prior value, so functionality of $T$ makes those new
values coextensional. The arbitrary graph objects occur only as
parameters of the already derived induction theorem. No graph
restriction, graph object coding, or filtering with a free graph
parameter is used.
\end{proof}

\begin{theorem}[Existence of bounded numerical histories]
\label{un:history-total}
For every numerical initial value $a$ and length $b$, there are numerical
$c,d$ and an object $f$ satisfying $H_T(f;a,b,c,d;\bar t)$.
\end{theorem}
\begin{proof}
Induct on $b$ with the numerical parameters $a,\bar t$ retained.
At an empty $b$, choose a number $n$ containing $a,b,\bar t$ and use
Theorem~\ref{un:graphs} with
\[
                         R(i,y):=(i\Eq b\land y\Eq a).
\]
This supplies a history on the single index class $[b]$, with terminal
value and bound $a$. No successor has an empty output, so there are no
transition conditions between indices in this domain.

For the step, let $\Succ(b,b')$. Take an old history with terminal
value $c$ and bound $d$, and choose $c'$ satisfying $T(c,c';\bar t)$.
Comparison supplies a numerical $d'$ at least $d$ and $c'$. Choose a
number $n$ containing the finite numerical tuple
$a,b,b',c,c',d,d',\bar t$. Apply Theorem~\ref{un:graphs} to
\begin{equation}\label{un:canonical-extension}
\begin{split}
 R(i,y;a,b,b',c,c',d,\bar t):={}&
 \bigl(i\preccurlyeq b\land
   \exists h\,[H_T(h;a,b,c,d;\bar t)\land V_h(i,y)]\bigr)\\
 &\lor(i\Eq b'\land y\Eq c').
\end{split}
\end{equation}
Every free parameter is numerical; the old graph variable $h$ is
bound. The formula is invariant in all its numerical arguments because
$H_T$ and $V_h$ are invariant, with the same witness $h$ after a change
of representatives.

There is at least one old history. By Lemma~\ref{un:history-unique},
all old histories in the first clause have the same values modulo
$\Eq$, and those values are at most $d$. Thus the first clause is
total and functional on $I_b$. The second clause adds the single
value class $[c']$ at the new index class $[b']$. The domains are
disjoint, and together they equal $I_{b'}$. All relevant indices and
values belong to $n$, so the numerical graph comprehension loses none.

The resulting graph retains the initial value and has terminal value
$c'$ and bound $d'$. An old transition agrees with the corresponding
transition in the original history by pointwise comparison and
invariance of $T$. For a transition into the new class $[b']$,
successor injectivity identifies its predecessor class as $[b]$.
Its value classes are $[c]$ and $[c']$, and the transition holds by
the choice of $c'$. This proves the induction step.
\end{proof}

Put
\[
 A_T(a,b,c;\bar t):=\exists f\exists d\,
                     [\N(d)\land H_T(f;a,b,c,d;\bar t)],
\]
with numerical guards on all displayed arithmetic arguments.

\begin{lemma}[Exact recursion equations]\label{un:recursion-equations}
The relation $A_T$ is total and functional modulo $\Eq$, and invariant
in all numerical arguments. It satisfies
\begin{align}
 \Emp(b)&\ \to\
       [A_T(a,b,c;\bar t)\leftrightarrow c\Eq a],
                                                     \label{un:rec-zero}\\
 \Succ(b,b')&\ \to\
 \bigl[A_T(a,b',c';\bar t)\leftrightarrow
 \exists c\,(\N(c)\land A_T(a,b,c;\bar t)\land T(c,c';\bar t))\bigr].
                                                     \label{un:rec-step}
\end{align}
\end{lemma}
\begin{proof}
Totality is Theorem~\ref{un:history-total}. Pointwise comparison at the
terminal index proves functionality, and the defining formulas prove
invariance. The initial and terminal clauses yield the forward zero
equation; the base construction, followed by invariance, yields its
reverse direction.

For the forward successor equation, take a history on $I_{b'}$ with
terminal value $c'$ and some numerical bound $d$. Its value at $b$
gives a number $c$, and its last transition gives $T(c,c';\bar t)$.
To obtain an actual history on $I_b$, choose a number containing
$a,b,b',c',d,\bar t$ and use numerical graph comprehension with
\begin{equation}\label{un:canonical-restriction}
 i\preccurlyeq b\ \land\
       \exists h\,[H_T(h;a,b',c',d;\bar t)\land V_h(i,y)].
\end{equation}
Again all free parameters are numerical. The given history witnesses
totality of this relation. Pointwise comparison makes it functional and
identifies all its values with the old values, which remain at most
$d$. Its graph is an $I_b$-history with terminal value $c$. This proves
$A_T(a,b,c;\bar t)$ without separating an old graph using that graph
as an illicit matrix parameter.

Conversely, given an old history and $T(c,c';\bar t)$, repeat the
canonical extension~\eqref{un:canonical-extension} with this chosen
$c'$. Its output supplies the desired new history.
\end{proof}

\subsection{The interpretation of PA}

\begin{theorem}[Arithmetic in the unary theory]\label{un:pa-lower}
The theory $\UU$ interprets full first-order Peano arithmetic.
\end{theorem}
\begin{proof}
Use $\N$ as the interpreted domain, $\Eq$ as equality, the empty class
as zero, and $\Succ$ as successor. The domain is nonempty.
Corollary~\ref{un:bootstrap} and Lemma~\ref{un:local} prove totality
and injectivity of successor, the omission of zero, and the existence
of a predecessor for every nonzero class.

First take $T(u,v):=\Succ(u,v)$ and define
\[
                         \operatorname{Add}(a,b,c):=A_T(a,b,c).
\]
This transition satisfies all hypotheses of the recursion construction.
Thus addition is total, invariant, and functional on numerical classes,
and~\eqref{un:rec-zero}--\eqref{un:rec-step} give
\[
                       [a]+0=[a],\qquad
                       [a]+S[b]=S([a]+[b]).
\]
Next, with numerical parameter $a$, use the already defined transition
$T_a(u,v):=\operatorname{Add}(u,a,v)$ and define
\[
 \operatorname{Mul}(a,b,c):=
             \exists e\,[\Emp(e)\land A_{T_a}(e,b,c;a)].
\]
Addition is already total, functional, and invariant, so this second
recursion is justified. All empty representatives are coextensional;
hence multiplication is total and functional and is independent of
the chosen initial empty representative. Its equations are
\[
                       [a]\cdot0=0,\qquad
                       [a]\cdot S[b]=[a]\cdot[b]+[a].
\]

Translate arithmetic terms by these total functional graphs, equality
by $\Eq$, and quantifiers by relativization to $\N$. Structural
induction on formulas proves invariance under replacing numerical
representatives. The finite arithmetic axioms follow from the
successor facts and recursion equations. Every translated induction
instance, with all its parameters, is an instance of
Theorem~\ref{un:induction}. Therefore every axiom of PA is valid under
this fixed definable interpretation.
\end{proof}

The interpretation uses only numerical parameters in its graph
comprehensions. Arbitrary graph objects occur as parameters of the
already proved induction theorem or as bound variables in
\eqref{un:canonical-extension} and~\eqref{un:canonical-restriction}.
The quotient language above describes a first-order interpretation;
it does not assert the existence of a set of all numerical classes.

\subsection{Uniformity needed for proof transport}\label{un:uniformity}

The proofs above provide uniform finite derivation templates.
Parameter internalization, the fixed-length least-tuple formula, the
failure predicate, and the unary matrix~\eqref{un:bnc-matrix} are formed
by substitution and fresh bound-variable renaming. For bounded
numerical comprehension, proofs of invariance and support are explicit
inputs; no algorithm decides whether an arbitrary formula has those
properties. Only finitely many such formulas are needed to establish
the fixed arithmetic interpretation. The graph and history definitions
are finite formulas formed before their existence proofs; they do not
contain a satisfaction predicate or a syntactic fixed point.

After the arithmetic interpretation has been fixed, translating a PA
induction instance requires only inserting its translated formula in
Theorem~\ref{un:induction}. The remaining arithmetic axioms have fixed
finite derivations. Translation of terms to operation graphs,
relativization, formula expansion, substitution, proof concatenation,
and the finite loops over displayed parameter tuples are primitive
recursive syntactic operations. Structural induction on formulas and
induction on coded derivations verify their correctness in PA. These
facts justify the usual PA-verifiable refutation transport; they do not
constitute a delivered proof-assistant certificate or an implemented
end-to-end proof translator.

